% !TeX root = main.tex
\chapter{Notation Reference}
\label{chap:notation}

This appendix summarises the main notation used throughout the dissertation. The goal is not to restate every paper-specific symbol, but to keep the shared notation for the synchronous, submodel, and buffered-asynchronous methods in one place.

\section*{General Federated Learning}
\vspace{-1em}
\begin{longtable}{@{} p{0.18\linewidth} p{0.78\linewidth} @{}}
$\mathcal{K}$ & Set of all available clients. \\
$\mathcal{C}_t$ & Client set selected in synchronous round $t$. \\
$t$ & Synchronous communication-round index. \\
$s$ & Buffered-asynchronous server-step index. \\
$T$ & Total number of synchronous rounds. \\
$S$ & Total number of asynchronous server steps. \\
$w^{(t)}$ & Global model after synchronous round $t$. \\
$w^{[s]}$ & Global model after asynchronous server step $s$. \\
$[w]_j$ & Coordinate \(j\) of model \(w\); bracket notation distinguishes parameter coordinates from client-indexed models such as \(\tilde{w}_k\). \\
$D_k$ & Local training dataset stored on client $k$. \\
$n_k$ & Number of training examples used by client $k$. \\
$p_k$ & Client weight in the global objective; FedAvg normalisation over the selected client set is written explicitly over $\mathcal{C}_t$ in the main text. \\
$F_k(w)$ & Local objective of client $k$. \\
$F(w)$ & Global objective, typically $F(w)=\sum_{k \in \mathcal{K}} p_k F_k(w)$. \\
$\eta_\ell$ & Client-side learning rate for local optimisation. \\
$\eta_s$ & Server-side learning rate used in aggregation updates. \\
$E$ & Number of local epochs per synchronous round. \\
$\tilde{w}_k^{(t+1)}$ & Local model returned by client $k$ after training from synchronous global model $w^{(t)}$. \\
\end{longtable}

\section*{Submodel Training}
\vspace{-1em}
\begin{longtable}{@{} p{0.18\linewidth} p{0.78\linewidth} @{}}
$r_k$ & Fixed model rate assigned to client $k$. \\
$I_k^{(t)}$, $I_k^{[s]}$ & Parameter subset visible to client $k$ in synchronous round $t$ or asynchronous dispatch step $s$. \\
$w_{k,0}^{(t)}$ & Client-visible model or submodel used by client $k$ before local training. \\
$\mathcal{U}_j^{(t)}$ & Set of participating clients whose extracted submodels include global parameter $j$ at round $t$. \\
$\Psi$, $\Psi_{\mathcal{M}}$ & Method-specific extraction policy that maps a global model, model rate, and round index to the client-visible submodel indices. \\
$s_j^{(t)}$ & Magnitude-based importance score assigned to parameter $j$ for the \texttt{FIARSE} sparse mask in round \(t\). \\
\end{longtable}

\section*{Buffered-Asynchronous Training}
\vspace{-1em}
\begin{longtable}{@{} p{0.18\linewidth} p{0.78\linewidth} @{}}
$M$ & Maximum number of concurrent in-flight client requests in buffered asynchronous execution. \\
$K$ & Buffer size used by the server before applying an asynchronous update. \\
$\mathcal{B}$ & Mutable set of accepted updates currently stored in the async buffer. \\
$\mathcal{B}_s$ & Accepted buffer whose aggregate is applied at asynchronous server step $s$. \\
$v_i$ & Model version originally sent to the client that produced async update $i$. \\
$\Delta_i$ & Delta associated with async update $i$, using the sign convention $\Delta_i=w^{[v_i]}-\tilde{w}_i$. \\
$[\Delta_i]_j$, $[\bar{\Delta}^{[s]}]_j$ & Coordinate \(j\) of an individual async delta or buffered aggregate delta. \\
$\tau_i$ & Staleness of async update $i$, measured in server steps since version $v_i$. \\
$q_i$ & Nonnegative staleness score assigned to buffered async update $i$ before normalisation. \\
$\alpha_i$ & Normalised buffer weight applied to buffered async update $i$, usually $\alpha_i=q_i/\sum_{u\in\mathcal{B}_s}q_u$. \\
$g(\tau_i)$ & Staleness-weighting function used to construct a buffered score from update staleness. \\
\end{longtable}

\section*{Method-Specific Parameters}
\vspace{-1em}
\begin{longtable}{@{} p{0.18\linewidth} p{0.78\linewidth} @{}}
$a$ & Polynomial staleness-decay exponent for the weighted \texttt{FedBuff} variant. \\
$c_j^{[s]}$ & Observed stale-weighted coverage mass for parameter $j$ in accepted buffer $\mathcal{B}_s$. \\
$\gamma$ & FedCover coverage-gain exponent. \\
$\kappa_{\max}$ & Maximum FedCover inverse-coverage gain. \\
$\kappa_j^{[s]}$ & Capped inverse-coverage gain applied to parameter $j$ at server step $s$. \\
\end{longtable}
